% !TEX root = ../hw02.tex
For $n\geq1$, define
\[
  m_n=\lfloor10^n\sqrt2\rfloor,\qquad
  a_n=\frac{m_n}{10^n},\qquad b_n=\frac{m_n+1}{10^n},
\]
and let $I_n=[a_n,b_n]\cap\Q$. Since $m_n\in\Z$, both endpoints
are rational. Thus $I_n$ is nonempty (it contains $a_n$), bounded,
and closed relative to $\Q$, being the intersection of a closed
subset of $\R$ with $\Q$.

Because $\sqrt2$ is irrational, the definition of the floor function
gives $m_n<10^n\sqrt2<m_n+1$. Consequently,
\[
  a_n<\sqrt2<b_n,\qquad b_n-a_n=10^{-n}.
\]
Figure~\ref{fig:hw02:rational-nested-intervals} illustrates these
rational intervals around the excluded point $\sqrt2$.
\begin{marginfigure}
  \centering
  % !TEX root = ../../problems/hw02.tex
% Schematic [a_n,b_n] intersect Q; the irrational limit is never included.
\begin{tikzpicture}[interval diagram]
  \draw[densely dashed,black!45] (1.8,0.55) -- (1.8,-2.4);
  \node[interval label,above=3pt] at (1.8,0.55) {$\sqrt2$};
  \foreach \n/\left/\right in {1/0/3.6,2/0.6/3,3/1.2/2.4} {
    \draw (\left,{1-\n}) -- (\right,{1-\n});
    \node[interval closed] at (\left,{1-\n}) {};
    \node[interval closed] at (\right,{1-\n}) {};
    \node[interval open] at (1.8,{1-\n}) {};
    \node[interval label,above=3pt] at (\left,{1-\n}) {$a_\n$};
    \node[interval label,above=3pt] at (\right,{1-\n}) {$b_\n$};
  }
  \node at (1.8,-2.9) {$\vdots$};
\end{tikzpicture}

  \caption{Schematic intervals $[a_n,b_n]\cap\Q$.
    Hollow markers indicate that $\sqrt2\notin\Q$;
    all endpoints are rational.}
  \label{fig:hw02:rational-nested-intervals}
\end{marginfigure}

Multiplying by $10$ gives $10m_n<10^{n+1}\sqrt2<10m_n+10$.
Taking floors therefore yields
\[
  10m_n\leq m_{n+1}\leq10m_n+9.
\]
It follows that
\[
  a_{n+1}=\frac{m_{n+1}}{10^{n+1}}\geq a_n,
  \qquad
  b_{n+1}=\frac{m_{n+1}+1}{10^{n+1}}\leq b_n.
\]
Thus $I_{n+1}\subseteq I_n$. In fact, the inclusion is strict because
the lengths of the corresponding real intervals decrease by a
factor of $10$. By the density of $\Q$ in $\R$, any nonempty
open portion removed contains a rational number, so at least one
point of $I_n$ does not belong to $I_{n+1}$.

The inequalities
\[
  0<\sqrt2-a_n<10^{-n},\qquad
  0<b_n-\sqrt2<10^{-n}
\]
show that $a_n\to\sqrt2$ and $b_n\to\sqrt2$. If
$x\in\bigcap_{n=1}^{\infty}I_n$, then $x\in\Q$ and
$a_n\leq x\leq b_n$ for every $n$. Since both endpoint sequences
converge to $\sqrt2$, the squeeze theorem applied to the constant
sequence with value $x$ gives $x=\sqrt2$, contradicting $x\in\Q$.\sidenote{
  \textit{Squeeze theorem.} If $u_n\leq v_n\leq w_n$ and
  $u_n,w_n\to\ell$, then $v_n\to\ell$.}
Therefore
\[
  \bigcap_{n=1}^{\infty}I_n=\varnothing.
\]

What fails is the \emph{completeness of $\Q$}. The corresponding
closed intervals in $\R$ have intersection $\{\sqrt2\}$, but this
point is absent from $\Q$.
